% 三视图(两幅)：正视图为带水平直径的圆，俯视图为矩形加半圆；12pt，.6pt灰度
% 点修：已为横向并排（正视图、俯视图），图名居中、基线对齐；投影边线/尺寸/虚实未动
\begin{tikzpicture}[line width=.6pt,font=\normalsize,x=1cm,y=1cm,>=Stealth]
% 正视图：圆心在(1.5,1.5)，半径1.5；水平直径线向右略伸出
\begin{scope}[shift={(0,0)}]
\draw (1.5,1.5) circle (1.5);
\draw[thin] (0,1.5) -- (3.35,1.5);
% 顶部宽度 2r
\draw[thin] (0,1.5) -- (0,3.6);
\draw[thin] (3.0,1.5) -- (3.0,3.6);
\draw[thin,<->] (0,3.35) -- (3.0,3.35);
\node[above] at (1.5,3.35) {$2r$};
% 右侧半径 r（中心线到底切线）
\draw[thin] (3.0,1.5) -- (3.75,1.5);
\draw[thin] (3.0,0) -- (3.75,0);
\draw[thin,<->] (3.55,1.5) -- (3.55,0);
\node[right] at (3.6,0.75) {$r$};
\node[below] at (1.65,-0.6) {正视图};
\end{scope}
% 俯视图：矩形宽2r高2r，上方半圆半径r；分界线向两侧略伸出
\begin{scope}[shift={(5.6,0)}]
\draw (0,0) -- (0,3.0) -- (3.0,3.0) -- (3.0,0);
\draw (0,0) -- (3.0,0);
\draw (0,3.0) arc (180:0:1.5);
\draw[thin] (-0.35,3.0) -- (3.35,3.0);
% 左侧半径 r（半圆高度）
\draw[thin] (-0.55,4.5) -- (1.5,4.5);
\draw[thin] (-0.35,4.5) -- (-0.35,3.0);
\draw[thin,<->] (-0.35,4.5) -- (-0.35,3.0);
\node[left] at (-0.4,3.75) {$r$};
% 右侧高度 2r（矩形高度）
\draw[thin] (3.0,3.0) -- (3.75,3.0);
\draw[thin] (3.0,0) -- (3.75,0);
\draw[thin,<->] (3.55,3.0) -- (3.55,0);
\node[right] at (3.6,1.5) {$2r$};
\node[below] at (1.5,-0.6) {俯视图};
\end{scope}
\end{tikzpicture}
